\section{CM reduction and consequences}
\label{sec:reduction}

The remaining step in the main proof is to express CM Hodge classes
as pullbacks of balanced Weil classes.  Together with CM transport
and Proposition~\ref{prop:weil}, this will prove
Theorem~\ref{thm:main}.  We then derive two specialization consequences.

\subsection{CM classes as pullbacks of Weil classes}

We give the CM reduction in the form needed here.  The stronger
reduction to classes on abelian varieties of split Weil type is due to
Deligne and Andr\'e \cite{Andre1992}; see the reconstruction in
\cite[Theorem~1 and Section~4]{Milne2020}. Only balanced signatures
and a compatible polarization are needed for our application, and
we prove this form directly using rational idempotents on powers.

\begin{lemma}[CM reduction]
\label{lem:cm-reduction}
Let $A/\Qbar$ be an abelian variety of CM type with good reduction at $w$,
and let $r>0$.  For every rational Hodge class
$\gamma\in H^{2r}(A(\C),\Q(r))$, there are finitely many abelian varieties
$B_i/\Qbar$ with good reduction at $w$, rational homomorphisms
$f_i\colon A\to B_i$, and rational Weil classes
$\eta_i\in H^{2r}(B_i(\C),\Q(r))$ such that
\begin{equation}
\label{eq:cm-pullback-decomposition}
                 \gamma=\sum_i f_i^*\eta_i.
\end{equation}
Each $B_i$ has an action of a CM field $K$ for which $H_1(B_i(\C),\Q)$
has rank $2r$ over $K$, the Hodge signatures are $(r,r)$ at every
embedding of $K$, and a polarization induces complex conjugation on $K$.
The equality specializes simultaneously in all the realizations used in
Theorem~\ref{thm:main}.
\end{lemma}

\begin{proof}
We first identify the eigenspace components that can support a rational
Hodge class.  For each such component we construct a balanced abelian
factor of a power of $A$, whose Weil classes pull back onto that
component.  Taking the sum of these rational pullback spaces will then
give the decomposition over $\Q$.

\medskip\noindent
\emph{Balanced eigenspace components.}
Set $V=H_1(A(\C),\Q)$.  We work throughout up to isogeny.  There is a
CM algebra $E'$, a product of CM fields, acting on $A$ such that $V$ is
free of rank one over $E'$.  To obtain it, decompose $A$ up to isogeny
into simple CM factors and take the product of their CM fields, with a
separate label for each repeated factor.  Transporting this action back
along an isogeny gives the asserted action on $A$ itself.

Choose a finite Galois CM field $K\subset\Qbar$ containing the images of
all the factors of $E'$, and put $S=\Hom_{\Q\text{-alg}}(E',K)$.  Then
\[
                V_K:=V\otimes_{\Q}K=\bigoplus_{\rho\in S}V_\rho,
                \qquad \dim_K V_\rho=1.
\]
The Hodge decomposition is diagonal for this decomposition after
extension to $\C$.  We suppress the common Tate twist in the following
linear-algebra calculations.  Alternating forms of degree $2r$ decompose
as
\[
 \bigwedge^{2r}V_K^*
     =\bigoplus_{\substack{\Delta\subset S\\ |\Delta|=2r}}L_\Delta,
 \qquad
 L_\Delta=
 \bigwedge^{2r}\left(\bigoplus_{\rho\in\Delta}V_\rho\right)^*,
\]
where each form on the right is extended by zero on the other summands.
Every $L_\Delta$ is one-dimensional and has a definite Hodge type.

Suppose the component of $\gamma$ in $L_\Delta$ is nonzero.  Since
$\gamma$ is rational, its component in $L_{\sigma\Delta}$ is the
semilinear conjugate of that component and is nonzero for every
$\sigma\in\Gal(K/\Q)$.  The Hodge condition therefore implies that
every $\sigma\Delta$ contains $r$ lines of each of the two Hodge types.
We call such a subset balanced under all conjugates.

\medskip\noindent
\emph{An abelian factor with balanced signatures.}
Fix one of these subsets $\Delta$.  Form the abelian variety
$A\otimes_{\Q}K$ in the category up to isogeny: it is a power of $A$
whose index space is the underlying rational vector space of $K$.
The regular action of $K$ on this index space is an action by rational
endomorphisms of the power.  Its rational homology is
\[
 T=V\otimes_{\Q}K,\qquad
 \dim_\Q T=[K:\Q]\dim_\Q V.
\]
In contrast with the coefficient space $V_K$, we regard $T$ here as a
$\Q$-vector space, with $K$ acting on its second factor.  The abelian
variety remains over $\Qbar$.  The actions on the two factors give
an action of the rational algebra
\[
 E'\otimes_{\Q}K\simeq\prod_{\rho\in S}K,
 \qquad e\otimes k\longmapsto\bigl(\rho(e)k\bigr)_{\rho\in S}.
\]
Let $e_\Delta$ be the idempotent with coordinates $1$ on $\Delta$ and
$0$ elsewhere.  It is a rational endomorphism of $A\otimes_{\Q}K$,
so its image is an abelian variety $B_\Delta$ up to isogeny.  In
particular no Galois invariance of $\Delta$ is required for this
construction.  Write
\[
                   U=H_1(B_\Delta(\C),\Q)=e_\Delta T.
\]
The index-field action of $K$ preserves $U$.

Now extend the coefficients of $U$ to $K$, and use the fixed inclusion
$K\subset\C$ to identify the embeddings of $K$ with $\Gal(K/\Q)$.
Evaluation of the index factor at $\tau$ identifies its eigenspace in
$T\otimes_\Q K$ with $V_K$.  In that eigenspace the idempotent retains
exactly the summands indexed by $\tau\Delta$.  Thus
\begin{equation}
\label{eq:cm-embedding-summand}
                  U_\tau=\bigoplus_{\rho\in\Delta}V_{\tau\rho},
\end{equation}
and the same description holds after extension to $\C$.  Explicitly,
evaluation of the idempotent at the pair of embeddings $(\xi,\tau)$ is $1$ exactly
when $\xi=\tau\rho$ for some $\rho\in\Delta$.  Thus
$\dim_K U_\tau=2r$ for every $\tau$.  Since $K$ is Galois and
$\Delta$ is balanced under all conjugates, each $U_\tau$ has $r$ lines
of each Hodge type.  Consequently $U$ has rank $2r$ over $K$,
$\dim B_\Delta=r[K:\Q]$, and the required Hodge signatures are
$(r,r)$.

\medskip\noindent
\emph{A compatible polarization.}
Let $\psi$ be a
rational polarization form on $V$, and equip the index space $K$ with
the positive definite rational symmetric form
\[
                       t(k,l)=\Tr_{K/\Q}(k\bar l).
\]
The form $\Psi=\psi\otimes t$ is a rational polarization form on the
power $A\otimes_{\Q}K$: its alternating and Hodge properties follow
from those of $\psi$, and its positivity follows by tensoring the
positive symmetric form associated to $\psi$ with $t$.  A positive
integral multiple gives an actual polarization.  For $a\in K$,
\[
                       t(ak,l)=t(k,\bar a l),
\]
so the Rosati involution of $\Psi$ induces conjugation on the
index-field action.  Restrict $\Psi$ to $B_\Delta$.  This is again a
polarization, and because $U$ is $K$-stable its adjoint identity still
induces conjugation on $K$.  This restriction does not require
$e_\Delta$ to be self-adjoint.

\medskip\noindent
\emph{Pullback and descent of the rational span.}
The factor $B_\Delta$ now has the field action, balanced signatures,
and polarization needed to define its Weil classes.  The rational
homomorphism $f_\Delta\colon A\to B_\Delta$ defined on
homology by
\[
                      v\longmapsto e_\Delta(v\otimes 1)
\]
has the needed effect on Weil classes.  After splitting coefficients,
its component at the identity embedding of $K$ is the projection
\[
                 \bigoplus_{\rho\in S}V_\rho
                    \longrightarrow
                 \bigoplus_{\rho\in\Delta}V_\rho.
\]
Thus pullback maps $\bigwedge^{2r}U_{\id}^*$ isomorphically onto
$L_\Delta$.  More generally the component indexed by $\tau$ pulls back
onto $L_{\tau\Delta}$.  The rational Weil-class space $W_K^r(B_\Delta)$
is characterized, with the twist restored, by
\[
 W_K^r(B_\Delta)\otimes_{\Q}K
       =\bigoplus_{\tau:K\hookrightarrow K}
                         \bigwedge^{2r}U_\tau^*(r).
\]
In particular $f_\Delta^*W_K^r(B_\Delta)$ is a rational subspace whose
extension to $K$ contains $L_\Delta(r)$.

Take the sum of these rational pullback spaces over the finitely many
subsets supporting $\gamma$.  Its scalar extension contains every
component of $\gamma$, and therefore contains $\gamma$ itself.  Since
extension of scalars is injective on the quotient by a rational
subspace, the sum already contains $\gamma$ over $\Q$.  This proves
\eqref{eq:cm-pullback-decomposition} with rational Weil classes.  Notice
that an individual embedding component need not descend to $\Q$;
only the entire Weil-class space and its pullback are used in this
descent argument.

\medskip\noindent
\emph{Good reduction and specialization.}
All the constructions use finitely many rational endomorphisms of
powers of $A$ and their images.  They therefore exist over $\Qbar$ and
descend, together with the maps, to a common number field after a finite
extension.  Choose that extension so that $A$ has good reduction at the
place specified by $w$.  For any $\ell\ne p$, the rational Tate module
of $B_\Delta$ is a Galois-stable direct summand of the rational Tate
module of the power of $A$.  It is unramified at this place.  The
N\'eron--Ogg--Shafarevich criterion gives good reduction of
$B_\Delta$ there; it also shows that the isogenies used in the
construction do not affect this conclusion
\cite[Chapter~IV, Theorem~3.5]{MilneAVNotes}.  Rational homomorphisms between
the resulting good-reduction varieties specialize after clearing
denominators.  The equality of rational Hodge classes
\eqref{eq:cm-pullback-decomposition} is respected by their compatible
realizations and by specialization, with the same rational coefficients
in every component.  This proves the last assertion.
\end{proof}

\begin{corollary}
\label{cor:cm-pairing}
The Lefschetz pairing property holds for the specialization of every
rational Hodge class on a CM abelian variety over $\Qbar$ with good
reduction at $w$.
\end{corollary}

\begin{proof}
In codimension zero the class is a rational multiple of the unit, so
the assertion follows from the compatibility of divisor intersection
numbers with the normalized traces.  In positive codimension use
Lemma~\ref{lem:cm-reduction}.  Proposition~\ref{prop:weil} gives the
Lefschetz pairing property for the specialization of each $\eta_i$.
Lemma~\ref{lem:pairing-functorial} preserves this property under
rational pullback and rational linear combinations.  Applying it to
\eqref{eq:cm-pullback-decomposition} proves the assertion.
\end{proof}

\subsection{Proof of the main theorem}

\begin{proof}[Proof of Theorem~\ref{thm:main}]
The case \(r=0\) follows from divisor intersection numbers as in the
preceding corollary. This includes \(\dim A=0\), where the abelian
variety is a point. If \(r=\dim A\), a rational top class is a
rational multiple of the point class; its normalized trace is
preserved by specialization, and the complementary product is \(1\).
We may therefore assume \(0<r<\dim A\).

Corollary~\ref{cor:cm-transport} gives a CM abelian variety $B/\Qbar$ with good
reduction at $w$, a rational Hodge class $\kappa$ on $B$, a
quasi-isogeny $a\colon A_0^n\to B_0$, and the first-factor inclusion
$j\colon A_0\to A_0^n$, such that
\[
                            \gamma_0=(aj)^*\kappa_0
\]
simultaneously in all the realizations under consideration.
Corollary~\ref{cor:cm-pairing} gives the Lefschetz pairing property for
$\kappa_0$, and Lemma~\ref{lem:pairing-functorial} gives it for
$\gamma_0$.  Hence, for each prescribed product of complementary
divisor classes on $A_0$, all the normalized traces in
Theorem~\ref{thm:main} equal one rational number.  This proves the
theorem.
\end{proof}

\subsection{Specialization in a fixed good theory}

We now use CM transport to place the specialization of every rational
Hodge class in one fixed theory of rational Tate classes.

Continue to fix \(\Qbar\subset\C\), \(p\), and \(w\), with
\(\F\), \(\Kzero\), and \(\Cw\) as in the introduction. Let
\(\mathcal S_0\) be Milne's class of \(\F\)-varieties whose connected
components are products of abelian varieties and projective spaces.
The CM case of Theorem~\ref{thm:main}, by linearity for complementary
Lefschetz classes, satisfies the hypothesis of
\cite[Theorem~4.5]{Milne2009}. Fix one good theory \(\mathcal R\) of
rational Tate classes on \(\mathcal S_0\) supplied by that theorem,
before varying any abelian variety or Hodge class.

For \(\rho\in\mathcal R^r(X)\), write \((\rho)_\ell\) for its
\(\ell\)-adic component when \(\ell\ne p\), and \((\rho)_p\) for
its native crystalline Tate component in
\[
 H^{2r}_{\mathrm{cris}}(X/W(\F))[1/p](r).
\]
We use three properties of the theory.  Its spaces are rational vector
spaces preserved by pullback, its realization maps are injective, and
specializations of CM Hodge classes belong to it.  These are the
properties in axioms \textup{(R1)}, \textup{(R4)}, and the goodness
axiom \textup{(R3)}, respectively
\cite[Definitions~2.1 and~3.1]{Milne2009}.  The following consequence
extends the last property to all rational Hodge classes, while keeping
this one theory fixed.

\begin{corollary}[Specialization of all rational Hodge classes]
\label{cor:all-hodge-specialization}
For every abelian variety \(A/\Qbar\) of dimension \(d\) with good
reduction \(A_0/\F\) at \(w\), every \(0\leq r\leq d\), and every
rational Hodge class
\[
 \gamma\in H^{2r}(A(\C),\Q(r)),
\]
there is a unique class \(\rho_\gamma\in\mathcal R^r(A_0)\) such that
\[
 (\rho_\gamma)_\ell=\gamma_{0,\ell}\quad(\ell\ne p),\qquad
 (\rho_\gamma)_p\otimes 1=\gamma_{0,p}.
\]
The prime-to-\(p\) equalities are in
\(H^{2r}_{\acute{e}t}(A_0,\Q_\ell(r))\); the crystalline equality is in
\[
 H^{2r}_{\mathrm{cris}}(A_0/W(\F))[1/p](r)
       \otimes_{\Kzero}\Cw.
\]
The specialization components are those of Theorem~\ref{thm:main}.
The assertion holds for every prime \(p\),
including \(p=2\), with the prime, place, and theory fixed before
\(A\) and \(\gamma\) vary.
\end{corollary}

\begin{proof}
In codimension zero, take the corresponding rational multiple of the
unit, which belongs to the theory by pullback from the point.  This
also covers $d=0$.  In top codimension $r=d>0$, take $q$ times the
point class, where $q\in\Q$ is the normalized trace of $\gamma$.
The point class belongs to $\mathcal R^d(A_0)$
\cite[\S2.3]{Milne2009}, and its normalized trace is one.  Compatibility
of traces with specialization gives the required components in these
cases.

Suppose now that $0<r<d$.  Choose the CM variety $B$, Hodge class
$\kappa$, and rational homomorphism
\[
 b=aj\in\Hom^0(A_0,B_0),\qquad
 \gamma_0=b^*\kappa_0
\]
from Corollary~\ref{cor:cm-transport}.  The equality holds in all
the stated realizations, with the usual absolute-Hodge specializations
of $\kappa$.  Axiom \textup{(R3)} gives one class
\(\rho_\kappa\in\mathcal R^r(B_0)\) with
\[
 (\rho_\kappa)_\ell=\kappa_{0,\ell}\quad(\ell\ne p),\qquad
 (\rho_\kappa)_p\otimes1=\kappa_{0,p}.
\]

Choose an integer \(N>0\) such that \(h=Nb\) is an actual
homomorphism \(A_0\to B_0\). Pullback by \(h\) preserves the theory
by axiom \textup{(R1)} of \cite[Definition~2.1]{Milne2009}. Since
\(\mathcal R^r(A_0)\) is a \(\Q\)-vector space, we may set
\[
 \rho_\gamma=N^{-2r}h^*\rho_\kappa\in\mathcal R^r(A_0).
\]
Functoriality of the component maps and
\eqref{eq:rational-pullback} give
\begin{align*}
 (\rho_\gamma)_\ell
   &=N^{-2r}h^*\kappa_{0,\ell}
     =b^*\kappa_{0,\ell}=\gamma_{0,\ell}
       &&(\ell\ne p),\\
 (\rho_\gamma)_p\otimes1
   &=N^{-2r}h^*\kappa_{0,p}
     =b^*\kappa_{0,p}=\gamma_{0,p}.
\end{align*}
The second line is taken after extension from $\Kzero$ to $\Cw$.
Pullback in \textup{(R1)} applies to any homomorphism $h$, and rational
scalar multiplication allows any denominator $N$.  Thus the argument
requires neither an isogeny nor separability, and remains valid when
$p\mid N$, including at $p=2$.

Finally, two classes with the required components have the same
\(\ell\)-adic component for any one \(\ell\ne p\). Injectivity of
that component map by \textup{(R4)} makes them equal. Thus the class
is unique in the fixed theory, independently of the auxiliary choices.
\end{proof}

This cohomological membership statement alone gives no algebraic cycle
and proves neither the Hodge nor the Tate conjecture.  The construction
takes place at the fixed
reduction place: its CM auxiliary may have dimension $n\dim A$, and
the map $b$ is between special fibers.  It supplies no same-dimensional
CM lift or map between the generic fibers.

\subsection{Algebraicity after good reduction}

The CM Hodge theorem \cite[Theorem~1.1]{OpenAICMHodge2026}, combined
with our CM transport, strengthens the preceding specialization
statement. It is an additional input to the following consequence,
not to the proof of Theorem~\ref{thm:main}.

\begin{corollary}[Algebraic specialization of rational Hodge classes]
\label{cor:algebraic-specialization}
Fix the prime and place as in Theorem~\ref{thm:main}. Let
\(A/\Qbar\) be an abelian variety of dimension \(d\) with good
reduction \(A_0/\F\) at \(w\), and let
\(\gamma\in H^{2r}(A(\C),\Q(r))\) be a rational Hodge class,
where \(0\leq r\leq d\). There is one rational algebraic cycle
\(z_\gamma\in\operatorname{CH}^r(A_0)_{\Q}\) such that
\[
 \operatorname{cl}_\ell(z_\gamma)=\gamma_{0,\ell}
       \quad(\ell\ne p),\qquad
 \operatorname{cl}_{\mathrm{cris}}(z_\gamma)\otimes1=\gamma_{0,p}.
\]
The crystalline equality is in
\[
 H^{2r}_{\mathrm{cris}}(A_0/W(\F))[1/p](r)
       \otimes_{\Kzero}\Cw.
\]
This holds for every prime \(p\), including \(p=2\).
Consequently, for every
\(y\in\operatorname{CH}^{d-r}(A_0)_{\Q}\), the traces of
\(\gamma_{0,\ell}\operatorname{cl}_\ell(y)\) for \(\ell\ne p\)
and of \(\gamma_{0,p}\operatorname{cl}_{\mathrm{cris}}(y)\)
are all the same rational number \(\deg(z_\gamma\cdot y)\).
\end{corollary}

\begin{proof}
In codimension zero or top codimension, use the corresponding rational
multiple of the unit or the point class, as in the preceding proof.
This includes dimension zero.  For $0<r<d$, use the CM variety $B$,
rational Hodge class $\kappa$, and map
\[
 b=aj\in\Hom^0(A_0,B_0),\qquad \gamma_0=b^*\kappa_0
\]
given by Corollary~\ref{cor:cm-transport}, with the equality in every
stated realization.

By \cite[Theorem~1.1]{OpenAICMHodge2026}, \(\kappa\) is the
Betti class of a rational algebraic cycle on \(B_{\C}\). A
representative can be chosen over \(\Qbar\): spread each of its
finitely many components in the Hilbert scheme of the fixed projective
variety \(B\), and choose a \(\Qbar\)-point on the same
geometrically connected component. The universal flat family preserves
the Betti cycle class. Thus there is
\(Z\in\operatorname{CH}^r(B)_{\Q}\) with class \(\kappa\).
Descend \(B\) and \(Z\) to a number field and extend the chosen
place so that \(B\) has an abelian scheme model. Specialization gives
\(Z_0\in\operatorname{CH}^r(B_0)_{\Q}\). Compatibility of
cycle classes with smooth proper specialization and the
de Rham--crystalline comparison identifies its classes with
\(\kappa_{0,\ell}\) and \(\kappa_{0,p}\), respectively.
In particular the same rational cycle \(Z_0\) gives every realization.

Choose \(N>0\) with \(h=Nb\) an actual homomorphism
\(A_0\to B_0\), and put
\[
 z_\gamma=N^{-2r}h^*Z_0\in\operatorname{CH}^r(A_0)_{\Q}.
\]
Here \(h^*\) is the Chow pullback between smooth varieties, defined
by the graph and refined intersection; it requires neither flatness
nor separability of \(h\). Functoriality of cycle classes and
\eqref{eq:rational-pullback} give the required equalities. Rational
coefficients permit this denominator even if \(p\mid N\); the
crystalline realization has already inverted \(p\). Finally,
compatibility of cycle classes with products and normalized traces
gives the asserted intersection number for every \(y\).
\end{proof}

This conclusion concerns only the special fiber at the fixed place.
It neither asserts the Hodge conjecture for arbitrary abelian varieties
in characteristic zero nor lifts \(z_\gamma\) to the original
generic fiber. The representing cycle is not asserted to be unique
modulo rational equivalence, and no comparison between different
reduction places is made.
